\subsubsection{Dark Dimension Gravitons as Dark Matter --- arXiv:2209.09249}

\textbf{Authors:} Eduardo Gonzalo, Miguel Montero, Georges Obied, Cumrun Vafa

\paragraph{主な主張}
micron-size Dark Dimension において、Standard Model brane から重力の普遍結合で不可避に生成される massive spin-2 KK graviton tower を dark matter とし、主に軽い KK graviton へ cascade decay して総質量密度をほぼ保ちながら現在 $m_{\rm DM}\sim 1$--$100\,\mathrm{keV}$ へ移ることで、正しい dark-matter abundance と cosmological coincidence を自然に再現できると主張する \cite{Gonzalo:2022DarkGravitons}。

\paragraph{新規性}
Dark Dimension の graviton KK tower を具体的な Dynamical Dark Matter realization として構成し、mesoscopic dimension に isometry がない場合の KK-momentum violation を $\delta n$ と $\beta$ で有効的に parameterize して intra-tower decay を cosmological evolution に組み込んだ点が特徴である。

\paragraph{位置付け}
Dark Dimension graviton dark matter と後続の decaying-dark-graviton phenomenology の出発点であり、brane backreaction / isometry breaking の観点では、仮定された KK-number violation と $\delta n,\beta$ の microscopic origin を幾何から導出する際の基準点となる。
